\section{Piedras, jugadas y prisioneros}
\label{sec:piedras}

Las piedras son los elementos fundamentales. En \pgfgoname, se proporcionan comandos y opciones para dibujarlas de manera flexible, permitiendo personalizar su apariencia y comportamiento según las necesidades del diagrama.

\subsection{Piedras}

\begin{command}{\stone\opt{\oarg{options}}\marg{intersection}}
	Este comando dibuja una piedra individual en la intersección especificada. El argumento obligatorio es la coordenada (ej. «C2»), mientras que las opciones opcionales permiten ajustar su estilo, como color, tamaño o marca. Si no se indica el jugador (|b|, |w| o |n|), la piedra es \emph{neutral}.
\end{command}


\begin{command}{\stones\opt{\oarg{options}}\marg{intersection \opt{, intersection, ...}}}
	Similar a |\stone|, pero permite colocar múltiples piedras con la sintaxis correspondiente (pág.~\pageref{subsec:sintaxis-de-coordenadas}). Si no se indica el jugador, las piedras también son \emph{neutrales}: para piedras negras hay que escribir |\stones[b]{...}|.
\end{command}

Las piedras, marcas y etiquetas que caen fuera de la región visible de un tablero parcial no se dibujan. Una coordenada que no existe en el tablero (por ejemplo, |E9| en un tablero de 5$\times$5) se ignora y produce una advertencia.

Es posible cambiar una opci\'on en todos los jugadores mediante |\pgfgoset{stone = {...}}|. Esto es equivalente a |\pgfgoset{players = {all = {stone = {...}}}}|.


\begin{key}{/pgfgo/stone/fill = \meta{options}}
	Controla las opciones de relleno de una piedra.
\end{key}


\begin{key}{/pgfgo/stone/fill/color = \meta{color}}	
	Define el color de relleno de una piedra.
	
	\begin{codeexample}[]
		\pgfgoset{goban = {5, partial = B2:D4}}
		
		\begin{goban}
			\stones[b]{C2,D3}
			\stone[w]{B4}
		\end{goban}
		
		\begin{goban}
			\stones[fill = {color = cyan!80!black}]{C2,D3}
			\stone[fill = {color = magenta!80!black}]{B4}
		\end{goban}
	\end{codeexample}
\end{key}

\begin{key}{/pgfgo/stone/fill/opacity = \meta{float} (initially 1)}
	
	Controla la opacidad del relleno de la piedra, desde 0 (transparente) hasta 1 (opaco).
	
	\begin{codeexample}[]
		\pgfgoset{goban = {5, partial = B2:D4}}
		
		\begin{goban}
			\stones[b]{C2,D3}
			\stone[w]{B4}
		\end{goban}
		
		\begin{goban}
			\stones[b, fill = {opacity = 0.8}]{C2,D3}
			\stone[w, fill = {opacity = 0.5}]{B4}
		\end{goban}
	\end{codeexample}
\end{key}


\begin{key}{/pgfgo/stone/line = \meta{option} (initially true)}
	Controla las opciones de la frontera de una piedra.
\end{key}

\begin{key}{/pgfgo/stone/line/color = \meta{color} (initially black)}
	
	Establece el color del borde de la piedra, que por defecto es negro para todas las piedras.
	
	\begin{codeexample}[]
		\pgfgoset{goban = {5, partial = B2:D4}}
		
		\begin{goban}
			\stones[b]{C2,D3}
			\stone[w]{B4}
		\end{goban}
		
		\begin{goban}
			\stones[b, line = {color = magenta}]{C2,D3}
			\stone[w, line = {color = cyan}]{B4}
		\end{goban}
	\end{codeexample}
\end{key}

\begin{key}{/pgfgo/stone/line/opacity = \meta{float} (initially 1)}
	
	Define la opacidad del borde de la piedra, permitiendo hacerlo más sutil o transparente.
	
	\begin{codeexample}[]
		\pgfgoset{goban = {5, partial = B2:D4}}
		
		\begin{goban}
			\stones[b]{C2,D3}
			\stone[w]{B4}
		\end{goban}
		
		\begin{goban}
			\stones[b, line = {opacity = 0.5}]{C2,D3}
			\stone[w, line = {opacity = 0.8}]{B4}
		\end{goban}
	\end{codeexample}
\end{key}

\begin{key}{/pgfgo/stone/line/width = \meta{dimension} (initially 0.7pt)}
	
	Controla el grosor del borde de la piedra.
	
	\begin{codeexample}[]
		\pgfgoset{goban = {5, partial = B2:D4}}
		
		\begin{goban}
			\stones[b]{C2,D3}
			\stone[w]{B4}
		\end{goban}
		
		\begin{goban}
			\stones[b]{C2,D3}
			\stone[w, line = {width = 2pt}]{B4}
		\end{goban}
	\end{codeexample}	
\end{key}

\begin{key}{/pgfgo/stone/radius = \meta{dimension} (initially 5.5pt)}
	
	Establece el tamaño de la piedra mediante su radio.
	
	\begin{codeexample}[]
		\pgfgoset{goban = {5, partial = B2:D4}}
		
		\begin{goban}
			\stones[b]{C2,D3}
			\stone[w]{B4}
		\end{goban}
		
		\begin{goban}
			\stones[b, radius = 0.15cm]{C2,D3}
			\stone[w, radius = 0.25cm]{B4}
		\end{goban}
	\end{codeexample}
\end{key}

\begin{key}{/pgfgo/stone/scale = \meta{float} (initially 1)}
	
	Define un factor de escala que afecta todos los elementos de la piedra (radio, borde, marcas), manteniendo sus proporciones.

	\begin{codeexample}[]
		\begin{goban}[5, partial = B2:D4]
			\stones[b, mark = circle]{C2}
			\stones[b, scale = 0.7, mark = circle]{D3}
			\stone[w, scale = 1.3, mark = triangle]{B4}
		\end{goban}
	\end{codeexample}
\end{key}

\begin{key}{/pgfgo/stone/black}
	
	Asigna a la piedra las características del primer jugador (negro por defecto), incluyendo color de relleno y borde. Sin |black| (o |b|), |\stone| y |\stones| dibujan piedras neutrales.
	
	\begin{codeexample}[]
		\pgfgoset{goban = {5, partial = B2:D4}}
		
		\begin{goban}
			\stones[b]{C2,D3}
			\stone[w]{B4}
		\end{goban}
		
		\begin{goban}
			\stones[b]{C2,D3}
			\stone[black]{B4}
		\end{goban}
	\end{codeexample}
\end{key}

\begin{key}{/pgfgo/stone/b}
	
	Abreviatura de «black», con idéntica funcionalidad.
\end{key}

\begin{key}{/pgfgo/stone/white}
	
	Asigna a la piedra las características del segundo jugador (blanco por defecto), con relleno blanco y borde negro.
	
	\begin{codeexample}[]
		\pgfgoset{goban = {5, partial = B2:D4}}
		
		\begin{goban}
			\stones[b]{C2,D3}
			\stone[white]{B4}
		\end{goban}
		
		\begin{goban}
			\stones[white]{C2,D3}
			\stone[white]{B4}
		\end{goban}
	\end{codeexample}
\end{key}

\begin{key}{/pgfgo/stone/w}
	
	Abreviatura de «white», con la misma funcionalidad.
\end{key}

\begin{key}{/pgfgo/stone/neutral}
	La piedra es del jugador neutral. Es el jugador por defecto de |\stone| y de |\stones|.

	\begin{codeexample}[]
		\begin{goban}[5, partial = B2:D4]
			\stones[b]{C2}
			\stones[w]{D3}
			\stone{B4}
			\stone[n]{C3}
		\end{goban}
	\end{codeexample}
\end{key}

\begin{key}{/pgfgo/stone/n}
	Abreviatura de «neutral».
\end{key}

\begin{key}{/pgfgo/stone/shadow = \meta{options} (default true)}
	Sombra de la piedra, dibujada debajo: |true|, |false|, |angle|, |offset|, |opacity|, |color|, |scale| e |image|. Sin |image| tiene la silueta de la piedra; con |image|, es una imagen propia (tal cual, o solo su forma con |silhouette|). Las opciones se describen en la referencia de jugadores (pág.~\pageref{subsec:jugadores}).

	\begin{codeexample}[]
		\pgfgoset{goban = {5, partial = B2:D4}}

		\begin{goban}
			\stones[b, shadow = {true, angle = -45, offset = 1.5pt}]{C2,D3}
			\stone[w, shadow = {true, angle = -45, offset = 1.5pt}]{B4}
		\end{goban}
	\end{codeexample}
\end{key}

\begin{key}{/pgfgo/stone/shine = \meta{options} (default true)}
	Brillo de la piedra, dibujado encima: |true|, |false|, |angle|, |opacity|, |color| e |image|. Las opciones se describen en la referencia de jugadores (pág.~\pageref{subsec:jugadores}).

	\begin{codeexample}[]
		\pgfgoset{goban = {5, partial = B2:D4}}

		\begin{goban}
			\stones[b, shine]{C2,D3}
			\stone[w, shine = {true, color = gray}]{B4}
		\end{goban}
	\end{codeexample}

	Con |image|, el brillo es una imagen con transparencia (PNG o PDF), centrada en la piedra (tal cual, o solo su forma, del color del brillo, con |silhouette|). Si la piedra también es una imagen, el brillo se dibuja encima y no gira con ella: las piedras pueden tener vetas giradas al azar (|rotate = random|) y la luz siempre del mismo lado.

	\begin{codeexample}[]
		\begin{goban}[5, partial = B2:D4]
			\stones[w, image = {file = pgfgo-yunzi-white-1.png, rotate = random},
				shine = {true, image = {file = shine.png}}]{B2, C2, D2, B3, C3, D3}
		\end{goban}
	\end{codeexample}
\end{key}


\label{subsec:piedrasimagen}\begin{key}{/pgfgo/stone/image = \meta{options}}
	Dibuja la piedra con una imagen (PNG, JPG o PDF) en lugar del círculo: no se dibujan el relleno ni el borde, y la imagen no tiene por qué ser redonda. Si el brillo (|shine|) está prendido, se dibuja encima de la imagen. La sombra (|shadow|) se dibuja con la forma de la imagen; para eso la imagen tiene que tener transparencia (un PNG con canal alfa o un PDF): la sombra de un JPG es un rectángulo. Un prisionero con imagen se ve translúcido, como cualquier prisionero.

	Los archivos se buscan como con |\includegraphics| (por ejemplo, en las carpetas de |\graphicspath|) y, si no están, entre las imágenes que trae \pgfgoname, en su carpeta |img| (las de los perfiles |shell| y |yunzi|, pág.~\pageref{subsec:perfilshell}). Cada archivo, con cada tamaño, se guarda una sola vez en el PDF, por más piedras que lo usen (pág.~\pageref{subsubsec:imagenespdf}). La sombra con la forma de la imagen necesita pdf\LaTeX\ o Lua\LaTeX\ (con salida PDF); con otros motores la piedra se dibuja sin sombra y con un aviso.

	Un nombre de archivo con coma va entre llaves dobles, |file = {{a,b.png}}|; con llaves simples, la coma separa dos archivos. Como con |\includegraphics|, el nombre no puede tener |&|, |#|, |%|, |$|, |{| ni |}|, ni dos espacios seguidos.

	Las imágenes SVG no se pueden usar directamente: hay que convertirlas antes a PDF, por ejemplo con |inkscape --export-type=pdf piedra.svg| o con |rsvg-convert -f pdf -o piedra.pdf piedra.svg|. Lo mismo vale para el fondo, la cuadrícula, los puntos de estrella y las marcas.

	\begin{codeexample}[]
		\pgfgoset{goban = {5, partial = B2:D4}}

		\begin{goban}
			\stones[b, image = {file = pgfgo-yunzi-black.png}, shadow = {true, angle = -45, offset = 1.5pt}]{C2,D3}
			\stone[w, image = {file = pgfgo-yunzi-white-1.png}, shadow = {true, angle = -45, offset = 1.5pt}]{B4}
		\end{goban}
	\end{codeexample}

	Para todas las piedras de un jugador, la imagen se da en sus opciones:

	\begin{codeexample}[]
		\pgfgoset{players = {
			black = {stone = {image = {file = pgfgo-yunzi-black.png}}},
			white = {stone = {image = {file = pgfgo-yunzi-white-1.png}}}}}

		\begin{goban}[5, partial = B2:D4]
			\moves[numbers = false]{C3, D3, C2, B3}
		\end{goban}
	\end{codeexample}
\end{key}

\begin{key}{/pgfgo/stone/image/file = \meta{file \opt{, file, ...}}}
	El archivo de la imagen. Con varios archivos, cada piedra usa uno, elegido por su posición: una piedra en un punto muestra siempre la misma imagen, en todos los diagramas. Sirve para variar el aspecto de las piedras, por ejemplo las vetas de las blancas.

	\begin{codeexample}[]
		\begin{goban}[5, partial = B2:D4]
			\stones[w, image = {file = {pgfgo-yunzi-white-1.png, pgfgo-yunzi-white-2.png}}]{B2, C2, D2, B3, C3, D3}
		\end{goban}
	\end{codeexample}
\end{key}

\begin{key}{/pgfgo/stone/image/width = \meta{dimension} (initially el diámetro de la piedra)}
	El ancho de la imagen; el alto mantiene la proporción. Por defecto es el diámetro de la piedra: $2 \times$ |radius| $\times$ |scale|.
\end{key}

\begin{key}{/pgfgo/stone/image/rotate = \meta{ángulo}, |random| o \meta{mínimo}|:|\meta{máximo} (initially 0)}
	Gira la imagen \meta{ángulo} grados, en sentido antihorario. Con |random|, cada piedra se gira un ángulo distinto, elegido por su posición: una piedra en un punto queda siempre igual, en todos los diagramas y en todos los cuadros de una animación. Sirve para que las vetas de una misma imagen no se vean repetidas. La sombra gira con la imagen, pero sigue cayendo hacia el mismo lado.

	\begin{codeexample}[]
		\begin{goban}[5, partial = B2:D4]
			\stones[w, image = {file = pgfgo-yunzi-white-1.png,
				rotate = random}]{B2, C2, D2, B3, C3, D3}
		\end{goban}
	\end{codeexample}

	Con un rango \meta{mínimo}|:|\meta{máximo}, el ángulo se elige igual que con |random|, pero entre \meta{mínimo} y \meta{máximo} grados (|random| es lo mismo que |0:360|). Sirve para imágenes que no se pueden girar del todo, como una piedra con luz de un lado.

	\begin{codeexample}[]
		\begin{goban}[5, partial = B2:D4]
			\stones[w, image = {file = pgfgo-yunzi-white-1.png,
				rotate = -20:20}]{B2, C2, D2, B3, C3, D3}
		\end{goban}
	\end{codeexample}
\end{key}

\begin{key}{/pgfgo/stone/image/opacity = \meta{float} (initially 1)}
	La opacidad de la imagen.

	\begin{codeexample}[]
		\begin{goban}[5, partial = B2:D4]
			\stones[b, image = {file = pgfgo-yunzi-black.png}]{C2}
			\stones[b, image = {file = pgfgo-yunzi-black.png, width = 8pt}]{D3}
			\stone[w, image = {file = pgfgo-yunzi-white-1.png, opacity = 0.5}]{B4}
		\end{goban}
	\end{codeexample}
\end{key}

\begin{stylekey}{/pgfgo/stone/image/false}
	Quita la imagen: la piedra vuelve a ser un círculo.
\end{stylekey}

\begin{key}{/pgfgo/stone/mark = \meta{mark keys}}
	
	Coloca una marca sobre la piedra (ej. triángulo, círculo). Las marcas pueden personalizarse ulteriormente con otras claves, como color o tamaño.
	
	\begin{codeexample}[]
		\pgfgoset{goban = {5, partial = B2:D4}}
		
		\begin{goban}
			\stones[b]{C2,D3}
			\stone[w]{B4}
		\end{goban}
		
		\begin{goban}
			\stones[b, mark = triangle]{C2,D3}
			\stone[w, mark = t-circle]{B4}
		\end{goban}
	\end{codeexample}
\end{key}

\begin{key}{/pgfgo/stone/forget}

	|resume| no trae esta piedra (p.~\pageref{subsec:recordar}). Queda en la memoria del |goban| con nombre: |\pgfgoview| la muestra con |forgotten|.

	\begin{codeexample}[]
		\pgfgoset{goban = {5, partial = B2:D4}}

		\begin{goban}[name = a]
			\stones[b]{C2,D3}
			\stone[w, forget]{B4}
		\end{goban}

		\begin{goban}[resume = a]
		\end{goban}
	\end{codeexample}
\end{key}

\begin{key}{/pgfgo/stone/prisoner=\meta{prisoner options} (default fill opacity = 0.5)}
	
	Caracteriza a una piedra como prisionera, típicamente con una opacidad reducida (0.5 por defecto) para indicar que ha sido capturada. Pueden añadirse otras opciones para personalizar su apariencia: se aplican después del estilo de los prisioneros (|/pgfgo/prisoner| y |/pgfgo/players/white/stone/prisoner|).
	
	\begin{codeexample}[]
		\pgfgoset{goban = {5, partial = B2:D4}}
		
		\begin{goban}
			\stones[b]{C2,D3}
			\stone[w]{B4}
		\end{goban}
		
		\begin{goban}
			\stones[b, prisoner]{C2,D3}
			\stone[white]{B4}
		\end{goban}
	\end{codeexample}
\end{key}

\subsubsection{Imágenes y tamaño del PDF}
\label{subsubsec:imagenespdf}

Con pdf\LaTeX\ o Lua\LaTeX\ (con salida PDF), cada imagen (de las piedras, el fondo, la cuadrícula, los hoshi, las marcas, el brillo o la sombra) se guarda en el PDF una sola vez, la primera vez que se usa; después, en todo el documento, cada piedra o baldosa que la usa solo la nombra. Una partida dibujada con |profile = shell| guarda la madera, las seis piedras y la sombra una vez, aunque tenga cientos de piedras repartidas en muchos diagramas. Girar una imagen (|rotate|) no agrega otra copia, y la sombra con la forma de la imagen (también con |silhouette|) usa la misma copia. Lo que sí agrega una copia es otro tamaño: la misma imagen con otro |width|, o en tableros con otro |scale| (por ejemplo, con |fit|), se guarda una vez por cada tamaño. Con otros motores (Xe\LaTeX, o salida DVI), la imagen se incluye cada vez que se usa, y el PDF crece con cada piedra.

Las piedras dibujadas, sin imagen, no se guardan así: cada una es un círculo y unas pocas instrucciones de dibujo, casi lo mismo que nombrar una copia guardada. Para piedras vectoriales más elaboradas (degradados, vetas dibujadas), conviene hacerlas en un PDF y usarlo con |stone/image|: un PDF es una imagen vectorial, se ve nítido con cualquier escala y se guarda una sola vez, como las demás imágenes.





\subsection{Movimientos}
\label{subsec:movimientos}

|\moves| representa secuencias de jugadas alternando automáticamente entre negro y blanco. Sus opciones deciden si las jugadas llevan número y personalizan las piedras asociadas.

\begin{command}{\moves\opt{\oarg{options}}\marg{intersection \opt{, intersection, ...}}}

	Dibuja piedras como parte de una secuencia de jugadas, alternando colores automáticamente (negro, blanco, negro, etc.) y con su número (|numbers|). Las jugadas son una secuencia: se juegan en el orden en que se escriben y una intersección puede repetirse (por ejemplo, para retomar un ko con |captures|). Acepta la sintaxis de |\stones| (pág.~\pageref{subsec:sintaxis-de-coordenadas}) y juega cada elemento en su orden (pág.~\pageref{subsec:orden-coordenadas}): un camino (|A1 -- A3 -- C3|) va de su primer punto al último, y un segmento, del primer extremo al segundo (|C:A 4| va de C4 a A4; |C5:1|, de C5 a C1). Lo que no tiene un recorrido fijo (rectángulos, |r|, columnas y filas enteras, |*|) se juega igual, con un aviso. Los grupos sustractivos (|;|) dan error: |\moves| solo agrega jugadas.

	\begin{codeexample}[]
		\begin{goban}[5]
			\moves{B1 -| D3, C:A 4, E5}
		\end{goban}
	\end{codeexample}
\end{command}

Dentro de un mismo |goban|, la numeración y el turno continúan de un |\moves| al siguiente. Cada |goban| nuevo empieza en la jugada 1 con negras; para continuar una partida en otro diagrama, se usan |from| y |b| o |w|, o bien |resume|.

Las opciones de |\moves| también aceptan las del estilo de la piedra (|fill|, |line|, |radius|, |shine|, |image|, etc.) y |mark|, como |\stones|. Por ejemplo, para marcar la última jugada:

\begin{codeexample}[]
	\begin{goban}[5, partial = B2:D4]
		\moves[numbers = false]{C3, D3, C2}
		\moves[numbers = false, mark = circle]{B3}
	\end{goban}
\end{codeexample}

Es posible cambiar las opciones de todos los jugadores mediante |\pgfgoset{move = {...}}| que es equivalente a |\pgfgoset{players = {all = {move = {...}}}}|, salvo |numbers|, que no es de un jugador: vale para todas las jugadas.

\begin{key}{/pgfgo/move/numbers=\meta{boolean} (default true, initially true)}

	Si las jugadas llevan su número. Sin números, las jugadas igual cuentan: la jugada siguiente sigue la numeración y el turno. Como opción de |\moves| vale para esas jugadas; con |\pgfgoset{move = {numbers = false}}| o |\begin{goban}[move = {numbers = false}]|, para todas las que siguen.

	\begin{codeexample}[]
		\begin{goban}[5, partial = B2:D4]
			\moves[numbers = false]{C3, D3}
			\moves{C2, B3}
		\end{goban}

		\begin{goban}[5, partial = B2:D4, move = {numbers = false}]
			\moves{C3, D3}
			\moves[numbers]{C2, B3}
		\end{goban}
	\end{codeexample}
\end{key}

\begin{key}{/pgfgo/move/black}
	
	Establece que la secuencia de movimientos comienza con el jugador negro, útil para forzar el inicio de una jugada en diagramas específicos.
	
	\begin{codeexample}[]
		\pgfgoset{goban = {5, partial = B2:D4}}
		
		\begin{goban}
			\stones[b]{C2,D3}
			\stone[w]{B4}
		\end{goban}
		
		\begin{goban}
			\moves[numbers = false, b]{C2,B4}
			\moves[numbers = false]{D3}
		\end{goban}
	\end{codeexample}
\end{key}

\begin{key}{/pgfgo/move/white}
	
	Establece que la secuencia comienza con el jugador blanco, alterando el orden natural de las jugadas.
	
	\begin{codeexample}[]
		\pgfgoset{goban = {5, partial = B2:D4}}
		
		\begin{goban}
			\stones[b]{C2,D3}
			\stone[w]{B4}
		\end{goban}
		
		\begin{goban}
			\moves[numbers = false, white]{C2,B4}
			\moves[numbers = false]{D3}
		\end{goban}
	\end{codeexample}
\end{key}

\begin{key}{/pgfgo/move/b}
	Abreviatura de |black|.
\end{key}

\begin{key}{/pgfgo/move/w}
	Abreviatura de |white|.
\end{key}

\begin{key}{/pgfgo/move/from = \meta{number}}
	
	Define el número inicial de una secuencia numerada de movimientos, permitiendo continuar desde un punto específico en lugar de empezar en 1.
	
	\begin{codeexample}[]
		\pgfgoset{goban = {5, partial = B2:D4}}
		
		\begin{goban}
			\stones[b]{C2,D3}
			\stone[w]{B4}
		\end{goban}
		
		\begin{goban}
			\moves[from = 7]{C2,B4}
			\moves{D3}
		\end{goban}
	\end{codeexample}
\end{key}

\begin{key}{/pgfgo/move/label = \meta{options}}
	
	Controla las propiedades de las etiquetas numéricas en los movimientos estrellados, como tamaño, color o formato.
\end{key}

\begin{key}{/pgfgo/move/label/font size=\meta{dimension} (initially automático)}

	Ajusta el tamaño de la fuente de las etiquetas, permitiendo adaptarlas al tamaño de las piedras o al diseño general. Si no se indica, el tamaño se ajusta solo según la cantidad de cifras del número.
	
	\begin{codeexample}[]
		\pgfgoset{goban = {5, partial = B2:D4}}
		
		\begin{goban}
			\stones[b]{C2,D3}
			\stone[w]{B4}
		\end{goban}
		
		\begin{goban}
			\moves[from = 7]{C2,B4}
			\moves[label = {font size = 0.5em}]{D3}
		\end{goban}
	\end{codeexample}
\end{key}

\begin{key}{/pgfgo/move/label/format = \meta{options} (initially arabic)}
	
	Define el estilo de numeración de las etiquetas.
\end{key}

\begin{key}{/pgfgo/move/label/format/arabic}
	
	Usa números arábigos, el formato por defecto para las etiquetas de movimientos.
	
	\begin{codeexample}[]
		\pgfgoset{goban = {5, partial = B2:D4}}
		
		\begin{goban}
			\stones[b]{C2,D3}
			\stone[w]{B4}
		\end{goban}
		
		\begin{goban}
			\moves[from = 7]{C2,B4}
			\moves{D3}
		\end{goban}
	\end{codeexample}
\end{key}

\begin{key}{/pgfgo/move/label/format/alph}
	
	Usa letras minúsculas para las etiquetas.
	
	\begin{codeexample}[]
		\pgfgoset{goban = {5, partial = B2:D4}}
		
		\begin{goban}
			\stones[b]{C2,D3}
			\stone[w]{B4}
		\end{goban}
		
		\pgfgoset{move={label = {format=alph}}}
		\begin{goban}
			\moves[from = 7]{C2,B4}
			\moves{D3}
		\end{goban}
	\end{codeexample}
\end{key}

\begin{key}{/pgfgo/move/label/format/Alph}
	
	Usa letras mayúsculas.
	
	\begin{codeexample}[]
		\pgfgoset{goban = {5, partial = B2:D4}}
		
		\begin{goban}
			\stones[b]{C2,D3}
			\stone[w]{B4}
		\end{goban}
		
		\pgfgoset{move = {label = {format = Alph}}}
		\begin{goban}
			\moves[from = 7]{C2,B4}
			\moves{D3}
		\end{goban}
	\end{codeexample}
\end{key}

\begin{key}{/pgfgo/move/label/format/roman}
	
	Usa números romanos minúsculos.
	
	\begin{codeexample}[]
		\pgfgoset{goban = {5, partial = B2:D4}}
		
		\begin{goban}
			\stones[b]{C2,D3}
			\stone[w]{B4}
		\end{goban}
		
		\pgfgoset{move = {label = {format = roman}}}
		\begin{goban}
			\moves{C2,B4}
			\moves{D3}
		\end{goban}
	\end{codeexample}
\end{key}

\begin{key}{/pgfgo/move/label/format/Roman}
	
	Usa números romanos mayúsculos.
	
	\begin{codeexample}[]
		\pgfgoset{goban = {5, partial = B2:D4}}
		
		\begin{goban}
			\stones[b]{C2,D3}
			\stone[w]{B4}
		\end{goban}
		
		\pgfgoset{move = {label = {format = Roman}}}
		\begin{goban}
			\moves{C2,B4}
			\moves{D3}
		\end{goban}
	\end{codeexample}
\end{key}

\begin{key}{/pgfgo/move/label/color = \meta{color}}
	
	Define el color de las etiquetas.
	
	\begin{codeexample}[]
		\pgfgoset{goban = {5, partial = B2:D4}}
		
		\begin{goban}
			\moves{C2,B4}
			\moves{D3}
		\end{goban}
		
		\begin{goban}
			\moves[label = {color = red}]{C2,B4}
			\moves[label = {color = blue}]{D3}
		\end{goban}
	\end{codeexample}
\end{key}

\begin{key}{/pgfgo/move/label/forget}

	Con |resume = {..., numbers}|, esta jugada viene sin su número (p.~\pageref{subsec:recordar}). Con |\pgfgoview|, el número solo se ve con |forgotten|.

	\begin{codeexample}[]
		\pgfgoset{goban = {5, partial = B2:D4}}

		\begin{goban}[name = a]
			\moves{C2}
			\moves[label = forget]{B4}
			\moves{D3}
		\end{goban}

		\begin{goban}[resume = {a, numbers}]
			\moves{C:D4}
		\end{goban}
	\end{codeexample}
\end{key}

\begin{key}{/pgfgo/move/forget}

	|resume| no trae estas jugadas, ni la piedra ni el número. La numeración sigue desde la última jugada que trae.
\end{key}


\subsection{Prisioneros}

\pgfgoname\ ofrece comandos específicos para representar prisioneros, generalmente con una opacidad reducida para distinguirlas de las piedras activas.

\begin{command}{\prisoner\opt{\oarg{options}}\marg{intersection}}
	
	Dibuja una piedra prisionera en una intersección específica, con opciones para personalizar su estilo.
\end{command}

\begin{command}{\prisoners\opt{\oarg{options}}\marg{intersection \opt{, intersection, ...}}}
	
	Coloca múltiples piedras prisioneras en una sola línea, eficiente para mostrar grupos capturados.
\end{command}

\begin{key}{/pgfgo/prisoner=\meta{stone options} (default fill opacity = 0.5)}
	
	Define las características de las piedras prisioneras, con una opacidad de relleno de 0.5 por defecto para indicar su estado. Pueden aplicarse otras opciones de |\stone|, como color o tamaño. Vale para los prisioneros de los tres jugadores; el estilo de cada uno (|/pgfgo/players/white/stone/prisoner|) va después.

	\begin{codeexample}[]
		\pgfgoset{goban = {5, partial = B2:D4}}

		\begin{goban}
			\stones[b]{C2}
			\prisoners[b]{D3}
			\prisoner[w]{B4}
		\end{goban}

		\pgfgoset{prisoner = {fill = {opacity = 0.2}, line = {color = red}}}

		\begin{goban}
			\stones[b]{C2}
			\prisoners[b]{D3}
			\prisoner[w]{B4}
		\end{goban}
	\end{codeexample}
\end{key}
